
\begin{exercisewrap}
\begin{nexercise} \label{croakerWhiteFishPacificExerConditions}
\index{data!ikan croaker putih dan merkuri|(}
Laman FDA menyediakan sejumlah data tentang kandungan merkuri pada ikan.
Berdasarkan sampel 15 ikan croaker putih (Pasifik),
rata-rata sampel dan simpangan baku masing-masing dihitung sebesar 0.287
dan 0.069 ppm (bagian per sejuta).
Ke-15 observasi tersebut berkisar dari 0.18 hingga 0.41 ppm.
Kita akan menganggap observasi-observasi ini saling independen.
Berdasarkan statistik ringkasan data,
adakah alasan untuk meragukan syarat normalitas
bagi observasi individual?\footnotemark{}
\end{nexercise}
\end{exercisewrap}
\footnotetext{Ukuran sampel kurang dari 30,
  jadi kita memeriksa pencilan yang jelas:
  karena semua observasi berada dalam 2 simpangan baku
  dari rata-rata, tidak ada pencilan yang jelas semacam itu.}

\begin{examplewrap}
\begin{nexample}{Estimasikan galat baku bagi
    $\bar{x} = 0.287$ ppm dengan menggunakan ringkasan data dalam
    Latihan Terarah~\ref{croakerWhiteFishPacificExerConditions}.
    Jika kita hendak menggunakan distribusi $t$ untuk menyusun
    interval kepercayaan 90\% bagi rata-rata sebenarnya dari
    kandungan merkuri, tentukan derajat kebebasan
    dan $t^{\star}_{df}$.}
  \label{croakerWhiteFishPacificExerSEDFTStar}%
  Galat baku: $SE = \frac{0.069}{\sqrt{15}} = 0.0178$.

  Derajat kebebasan: $df = n - 1 = 14$.

  Karena tujuannya adalah interval kepercayaan 90\%,
  kita memilih $t_{14}^{\star}$ sehingga luas kedua ekor
  adalah 0.1:
  $t^{\star}_{14} = 1.76$.
\end{nexample}
\end{examplewrap}

\begin{onebox}{Interval kepercayaan bagi satu rata-rata}
  Setelah menentukan bahwa interval kepercayaan bagi satu rata-rata
  akan berguna dalam suatu penerapan,
  terdapat empat langkah untuk menyusun interval tersebut:
  \begin{description}
  \item[Persiapkan.]
      Tentukan $\bar{x}$, $s$, $n$, dan pilih
      tingkat kepercayaan yang ingin Anda gunakan.
  \item[Periksa.]
      Verifikasi syarat-syarat untuk memastikan bahwa $\bar{x}$
      mendekati normal.
  \item[Hitung.]
      Jika syarat-syarat terpenuhi, hitung $SE$,
      cari $t_{df}^{\star}$, dan susun intervalnya.
  \item[Simpulkan.]
      Tafsirkan interval kepercayaan dalam konteks
      masalah tersebut.
  \end{description}
\end{onebox}

\begin{exercisewrap}
\begin{nexercise}
\label{croakerWhiteFish90ci}
Dengan menggunakan informasi dan hasil Latihan Terarah~\ref{croakerWhiteFishPacificExerConditions} serta Contoh~\ref{croakerWhiteFishPacificExerSEDFTStar}, hitung interval kepercayaan 90\% bagi rata-rata kandungan merkuri pada ikan croaker putih (Pasifik).\footnotemark{}
\end{nexercise}
\end{exercisewrap}
\footnotetext{
  $\bar{x} \ \pm\ t^{\star}_{14} \times SE
      \ \to\  0.287 \ \pm\  1.76 \times 0.0178
      \ \to\ (0.256, 0.318)$.
  Kita yakin 90\% bahwa rata-rata kandungan merkuri
  pada ikan croaker putih (Pasifik) berada antara 0.256 dan 0.318 ppm.}

\begin{exercisewrap}
\begin{nexercise}
Interval kepercayaan 90\% dari
Latihan Terarah~\ref{croakerWhiteFish90ci}
adalah 0.256 ppm hingga 0.318 ppm.
Dapatkah kita mengatakan bahwa 90\% dari ikan croaker putih (Pasifik)
memiliki kadar merkuri antara 0.256 dan 0.318 ppm?\footnotemark{}
\end{nexercise}
\end{exercisewrap}
\footnotetext{
  Tidak. Interval kepercayaan hanya memberikan rentang
  nilai-nilai yang masuk akal bagi suatu parameter populasi,
  dalam hal ini rata-rata populasi.
  Interval itu tidak menjelaskan apa yang mungkin kita amati
  pada observasi individual.}

\index{data!ikan croaker putih dan merkuri|)}

%Now that we've whet \Comment{spelling?} your palette with confidence
%intervals for a mean, let's speed on through to
%hypothesis tests for the mean.


\subsection[Uji $t$ satu sampel]
    {Uji $\pmb{t}$ satu sampel}
\label{oneSampleTTests}
 
